\documentclass[a4paper,12pt]{article}
\usepackage[spanish]{babel}
\usepackage[utf8]{inputenc}
\usepackage[T1]{fontenc}
\usepackage[margin=1in]{geometry}
\usepackage{fancyhdr}
\usepackage{graphicx}
\usepackage{amsmath, amssymb}
\usepackage{booktabs}
\usepackage{float}
\usepackage{xcolor}
\usepackage{listings}
\usepackage{hyperref}
\usepackage{tikz}
\usepackage{enumitem}
\usepackage{tabularx}

\usetikzlibrary{shapes.geometric}

% ── Colores del sistema ───────────────────────────────────────────
\definecolor{navy}{RGB}{15,42,76}
\definecolor{gold}{RGB}{251,191,36}

\hypersetup{
  colorlinks=true,
  linkcolor=navy,
  urlcolor=navy,
  citecolor=navy,
}

% Spanish babel rotula las tablas como "Cuadro"; lo forzamos a "Tabla".
\renewcommand{\tablename}{Tabla}

% ── Encabezado y pie ──────────────────────────────────────────────
\setlength{\headheight}{15pt}
\pagestyle{fancy}
\fancyhf{}
\fancyhead[L]{\small\textbf{FINESI}}
\fancyhead[C]{\small Universidad Nacional del Altiplano}
\fancyhead[R]{\small\today}
\fancyfoot[C]{\small Página \thepage}
\renewcommand{\headrulewidth}{0.4pt}
\renewcommand{\footrulewidth}{0.4pt}

% ── Estilo de tablas de casos de uso ──────────────────────────────
\newcolumntype{C}{>{\bfseries\raggedright\arraybackslash}p{3.4cm}}

% Figura de palo (stickman) reutilizable para el diagrama de CU.
% Uso: \stickman{x}{y}{Etiqueta}
\newcommand{\stickman}[3]{%
  \begin{scope}[shift={(#1,#2)}]
    \draw[thick,navy] (0,0.55) circle (0.18);      % cabeza
    \draw[thick,navy] (0,0.37) -- (0,-0.25);       % torso
    \draw[thick,navy] (-0.25,0.12) -- (0.25,0.12); % brazos
    \draw[thick,navy] (0,-0.25) -- (-0.20,-0.65);  % pierna izq
    \draw[thick,navy] (0,-0.25) -- (0.20,-0.65);   % pierna der
    \node[align=center, font=\footnotesize\bfseries, text=navy]
      at (0,-1.0) {#3};
  \end{scope}
}

\begin{document}

% ════════════════════════════════════════════════════════════════
%                              PORTADA
% ════════════════════════════════════════════════════════════════
\begin{titlepage}
\centering
\thispagestyle{empty}

\IfFileExists{logoU.png}{%
  \includegraphics[width=3cm]{logoU.png}\par\vspace{0.6cm}
}{%
  \fbox{\parbox[c][3cm][c]{3cm}{\centering\footnotesize logoU.png}}\par\vspace{0.6cm}
}

{\large\scshape Universidad Nacional del Altiplano -- Puno\par}
\vspace{0.3cm}
{\large Facultad de Ingeniería Estadística e Informática\par}
\vspace{0.15cm}
{\large Escuela Profesional de Ingeniería Estadística e Informática\par}

\vspace{2.2cm}
{\color{navy}\rule{\linewidth}{1.2pt}}\par
\vspace{0.8cm}
{\Huge\bfseries\color{navy} YachayQR\par}
\vspace{0.4cm}
{\LARGE Sistema SaaS de Control de Asistencia Escolar\par}
\vspace{0.8cm}
{\color{gold}\rule{\linewidth}{1.2pt}}\par

\vspace{1.0cm}
{\large\itshape Documentación del Proceso\par}
{\large\itshape Ingeniería de Software II\par}

\vspace{2.5cm}
\begin{flushleft}
\large
\textbf{Presentado por:}\hspace{0.3cm} Cliver Oliver Turpo Benique\par
\vspace{0.4cm}
\textbf{Curso:}\hspace{0.3cm} Ingeniería de Software II\par
\vspace{0.4cm}
\textbf{Docente:}\hspace{0.3cm} Ramiro Pedro Laura Murillo\par
\end{flushleft}

\vfill
{\large Puno -- Perú\par}
{\large \today\par}
\end{titlepage}

\tableofcontents
\newpage

% ════════════════════════════════════════════════════════════════
\section{Descripción del Sistema}
% ════════════════════════════════════════════════════════════════

\textbf{YachayQR} es un sistema de tipo SaaS (\textit{Software as a Service})
para el control de asistencia escolar mediante el escaneo del DNI del alumno,
que funciona simultáneamente como código de barras y código QR. El sistema
resuelve un problema cotidiano de los colegios: el \textbf{pase de lista
manual en papel}, un proceso lento, propenso a errores y que no deja registro
verificable de la hora exacta de ingreso ni permite avisar a los apoderados de
forma inmediata.

Cuando un alumno pasa su DNI por la lectora, el sistema lo identifica al
instante y registra automáticamente su estado del día como \textit{Presente}
o \textit{Tardanza} según la hora configurada en el horario escolar. Al cerrar
la jornada, los alumnos sin registro quedan marcados como \textit{Ausente} y
se notifica a sus apoderados por WhatsApp. Además gestiona justificaciones de
inasistencia con un límite controlado, detecta intentos de fraude (DNI
digitado a mano en lugar de escaneado) y genera reportes y carnets en PDF.

YachayQR es \textbf{multi-colegio}: cada institución funciona como un
inquilino (\textit{tenant}) totalmente aislado, con su propio subdominio y su
propia base de datos lógica. El sistema se encuentra \textbf{desplegado en
producción} en \url{https://yachayqr.com}, y cuenta con un colegio de
demostración con datos reales en \url{https://demo.yachayqr.com}.

\begin{figure}[H]
\centering
\begin{minipage}{0.45\linewidth}
  \centering
  \IfFileExists{captura_landing.png}{%
    \includegraphics[width=\linewidth]{captura_landing.png}
  }{%
    \fbox{\parbox[c][3cm][c]{\linewidth}{\centering
      \footnotesize \textbf{captura\_landing.png}\\[0.2cm]
      \url{https://yachayqr.com}}}
  }
  \caption{Landing page.}
\end{minipage}
\hfill
\begin{minipage}{0.45\linewidth}
  \centering
  \IfFileExists{captura_login.png}{%
    \includegraphics[width=\linewidth]{captura_login.png}
  }{%
    \fbox{\parbox[c][3cm][c]{\linewidth}{\centering
      \footnotesize \textbf{captura\_login.png}\\[0.2cm]
      \url{https://demo.yachayqr.com/login}\\[0.25cm]
      Usuario: \texttt{director}\\
      Contraseña: \texttt{demo1234}}}
  }
  \caption{Login del colegio demo (usuario: \texttt{director}, contraseña: \texttt{demo1234}).}
\end{minipage}
\end{figure}
\newpage

% ════════════════════════════════════════════════════════════════
\section{Entrevistas}
% ════════════════════════════════════════════════════════════════

Para identificar los requisitos del sistema se realizaron entrevistas con los
actores directamente involucrados en el control de asistencia de la
Institución Educativa Secundaria Túpac Amaru de Coasa (IESTA-COASA), provincia
de Carabaya, región Puno. A continuación se presentan las dos entrevistas más
representativas y los hallazgos que orientaron el diseño del sistema.

% ── Entrevista 1 ──────────────────────────────────────────────────
\subsection{Entrevista 1 -- Director de la institución}

\noindent
\begin{tabularx}{\linewidth}{C X}
\toprule
Fecha          & Mayo de 2026 \\
Entrevistador  & Cliver Oliver Turpo Benique \\
Entrevistado   & Prof. Hernán Chiapana Condori, Director de la I.E.S. Túpac Amaru de Coasa (IESTA-COASA), Carabaya -- Puno \\
Modalidad      & Entrevista semiestructurada presencial \\
\bottomrule
\end{tabularx}

\vspace{0.4cm}
\noindent\textbf{Preguntas y respuestas}

\begin{enumerate}[leftmargin=*, label=\textbf{P\arabic*.}]
  \item \textit{¿Cómo registran actualmente la asistencia de los alumnos?}\\
  ``Cada auxiliar pasa lista en la formación de la mañana con un cuaderno. Lo
  que anota lo pasa después a un cuaderno general, y a veces a un Excel. Es
  trabajo doble y se pierde tiempo valioso.''

  \item \textit{¿Qué problemas le genera ese proceso manual?}\\
  ``Los cuadernos se llenan de borrones y correcciones. Cuando un padre viene
  a reclamar que su hijo sí asistió, a veces no podemos demostrar la hora
  exacta en que entró. No hay una prueba confiable.''

  \item \textit{¿Cómo manejan las llegadas tarde?}\\
  ``Ese es el mayor dolor de cabeza. Marcar tardanzas a mano es subjetivo,
  depende de qué tan rápido escriba el auxiliar. No tenemos una hora de corte
  que el sistema aplique igual para todos.''

  \item \textit{¿Cómo se comunican con los padres cuando un alumno falta?}\\
  ``Casi no nos comunicamos el mismo día. El padre se entera al final de la
  semana o cuando lo citamos. Para entonces el alumno ya faltó tres veces y es
  tarde para actuar.''

  \item \textit{Si pudiera pedir una herramienta ideal, ¿qué debería hacer?}\\
  ``Que registre la entrada con el DNI, que avise al padre el mismo día por
  WhatsApp, y que yo pueda sacar un reporte por sección sin estar sumando a
  mano. Y que cada colegio tenga lo suyo, separado.''
\end{enumerate}

\vspace{0.2cm}
\noindent\textbf{Hallazgos clave}

\begin{table}[H]
\centering
\begin{tabularx}{\linewidth}{p{6.2cm} X}
\toprule
\textbf{Hallazgo} & \textbf{Cómo lo resuelve YachayQR} \\
\midrule
Doble registro manual (cuaderno + Excel) que consume tiempo.
& Registro único y automático al escanear el DNI. \\
\addlinespace
Falta de prueba confiable de la hora de ingreso.
& Cada asistencia guarda la hora exacta y el método de registro. \\
\addlinespace
Criterio subjetivo para marcar tardanzas.
& Hora de corte configurable; el sistema aplica Presente/Tardanza de forma uniforme. \\
\addlinespace
Comunicación tardía con los padres.
& Notificación por WhatsApp al apoderado el mismo día del cierre. \\
\addlinespace
Necesidad de reportes por sección y aislamiento por colegio.
& Reportes filtrables en Excel/PDF y arquitectura multi-colegio aislada. \\
\bottomrule
\end{tabularx}
\caption{Hallazgos de la entrevista con el Director.}
\end{table}

% ── Entrevista 2 ──────────────────────────────────────────────────
\subsection{Entrevista 2 -- Auxiliar de educación}

\noindent
\begin{tabularx}{\linewidth}{C X}
\toprule
Fecha          & Mayo de 2026 \\
Entrevistador  & Cliver Oliver Turpo Benique \\
Entrevistado   & Sra. Rosa Ccama Apaza, Auxiliar de educación del mismo colegio (nombre ficticio) \\
Modalidad      & Entrevista semiestructurada presencial \\
\bottomrule
\end{tabularx}

\vspace{0.4cm}
\noindent\textbf{Preguntas y respuestas}

\begin{enumerate}[leftmargin=*, label=\textbf{P\arabic*.}]
  \item \textit{¿En qué consiste su trabajo de registro cada mañana?}\\
  ``Me paro en la puerta con la lista impresa y voy marcando uno por uno
  mientras entran. Con 200 alumnos entrando casi al mismo tiempo, es
  imposible no equivocarse.''

  \item \textit{¿Qué tipo de errores ocurren con más frecuencia?}\\
  ``Marcar a un alumno que no vino, o no marcar a uno que sí entró. También
  confundir nombres parecidos. Después, pasar todo eso al cuaderno general es
  donde más se cometen errores.''

  \item \textit{¿Cuánto tiempo le toma todo el proceso?}\\
  ``El pase de lista en la formación son unos 20 minutos, pero pasar los datos
  en limpio y revisar me toma como una hora más cada día.''

  \item \textit{¿Qué pasa cuando un alumno llega tarde después de la formación?}\\
  ``Tengo que buscarlo en la lista y anotarlo aparte. Muchas veces esas
  tardanzas se me pasan o las anoto mal.''

  \item \textit{¿Le preocupa que alguien marque por otro compañero?}\\
  ``Sí, pasa. A veces un alumno me dice un nombre que no es el suyo para que
  no figure su tardanza. Con la lista a mano no tengo cómo comprobarlo.''
\end{enumerate}

\vspace{0.2cm}
\noindent\textbf{Hallazgos clave}

\begin{table}[H]
\centering
\begin{tabularx}{\linewidth}{p{6.2cm} X}
\toprule
\textbf{Hallazgo} & \textbf{Cómo lo resuelve YachayQR} \\
\midrule
Errores al marcar manualmente en hora pico de ingreso.
& El escaneo identifica al alumno por su DNI, sin elección manual. \\
\addlinespace
Confusión entre alumnos de nombres parecidos.
& La pantalla muestra foto, nombre completo y grado al escanear. \\
\addlinespace
Una hora diaria adicional pasando datos en limpio.
& No hay transcripción: el dato queda registrado y disponible al instante. \\
\addlinespace
Tardanzas posteriores a la formación que se omiten.
& El escaneo está disponible durante toda la ventana de horario. \\
\addlinespace
Riesgo de que un alumno se haga pasar por otro.
& Detección anti-fraude del ingreso manual y alerta al Director. \\
\bottomrule
\end{tabularx}
\caption{Hallazgos de la entrevista con la Auxiliar.}
\end{table}

% ════════════════════════════════════════════════════════════════
\section{Casos de Uso}
% ════════════════════════════════════════════════════════════════

A continuación se describen los casos de uso principales del sistema,
elaborados a partir de la funcionalidad realmente implementada en el código
fuente.

% ── CU01 ──────────────────────────────────────────────────────────
\subsection{CU01 -- Registrar asistencia por escaneo de DNI}

\begin{table}[H]
\centering
\begin{tabularx}{\linewidth}{C X}
\toprule
\textbf{Campo} & \textbf{Descripción} \\
\midrule
Actor & Auxiliar / Escáner (dispositivo físico). \\
\addlinespace
Precondición &
  Existe un horario escolar activo configurado; el usuario está autenticado y
  activo; el día actual es laborable y la hora está dentro de la ventana
  permitida. \\
\addlinespace
Flujo principal &
  1) El actor pasa el DNI del alumno por la lectora. \newline
  2) El sistema verifica que haya una sesión del día abierta; si es el primer
     escaneo del día, la abre automáticamente. \newline
  3) El sistema busca al alumno por su código de barras (su DNI). \newline
  4) Determina el estado: \textit{Presente} si la hora es anterior o igual a
     la hora de corte, \textit{Tardanza} si es posterior. \newline
  5) Registra la asistencia con la hora exacta y muestra en pantalla la foto,
     el nombre, el grado y el estado del alumno. \\
\addlinespace
Flujos alternativos &
  \textbf{4a. Alumno no encontrado:} el código no corresponde a ningún alumno
  activo $\rightarrow$ el sistema muestra ``Alumno no encontrado''. \newline
  \textbf{2a. Sesión ya cerrada:} la jornada del día fue cerrada $\rightarrow$
  el sistema rechaza el escaneo. \newline
  \textbf{2b. Fuera de horario:} la hora actual es anterior a la entrada o
  posterior al cierre $\rightarrow$ se informa el horario válido. \newline
  \textbf{2c. Día no laborable:} hoy no figura entre los días laborables
  $\rightarrow$ se rechaza el escaneo. \newline
  \textbf{5a. Doble registro:} el alumno ya fue registrado hoy $\rightarrow$ el
  sistema lo informa sin duplicar el registro. \\
\addlinespace
Postcondición &
  La asistencia del alumno queda registrada para la fecha; si era el primer
  escaneo, la sesión diaria queda abierta y los contadores actualizados. \\
\bottomrule
\end{tabularx}
\end{table}

% ── CU02 ──────────────────────────────────────────────────────────
\subsection{CU02 -- Cerrar sesión y notificar a los apoderados por WhatsApp}

\begin{table}[H]
\centering
\begin{tabularx}{\linewidth}{C X}
\toprule
\textbf{Campo} & \textbf{Descripción} \\
\midrule
Actor & Director / Auxiliar. \\
\addlinespace
Precondición & Existe una sesión del día abierta. \\
\addlinespace
Flujo principal &
  1) El actor solicita cerrar la jornada del día. \newline
  2) El sistema marca como \textit{Ausente} a todos los alumnos activos sin
     registro. \newline
  3) Actualiza los contadores de la sesión (presentes, tardanzas, ausentes,
     justificados). \newline
  4) Cambia el estado de la sesión a \textit{Cerrada}. \newline
  5) Encola una tarea en segundo plano (Celery) que envía mensajes de WhatsApp
     a los apoderados de los alumnos ausentes y con tardanza. \\
\addlinespace
Flujos alternativos &
  \textbf{1a. Sesión ya cerrada:} el sistema informa que la jornada ya fue
  cerrada y no realiza acciones. \newline
  \textbf{5a. Colegio sin plan WhatsApp activo:} la tarea se ejecuta pero omite
  el envío de mensajes en silencio (la notificación es parte del Plan
  Premium). \\
\addlinespace
Postcondición &
  La sesión queda cerrada, todos los alumnos tienen un estado definido y, si el
  plan está activo, los apoderados han sido notificados. \\
\bottomrule
\end{tabularx}
\end{table}

% ── CU03 ──────────────────────────────────────────────────────────
\subsection{CU03 -- Justificar inasistencia}

\begin{table}[H]
\centering
\begin{tabularx}{\linewidth}{C X}
\toprule
\textbf{Campo} & \textbf{Descripción} \\
\midrule
Actor & Psicólogo / Director (también el Auxiliar puede justificar). \\
\addlinespace
Precondición &
  Existe un registro de asistencia (tardanza o ausencia) que se desea
  justificar; el usuario tiene permiso para justificar. \\
\addlinespace
Flujo principal &
  1) El actor selecciona la asistencia del alumno a justificar. \newline
  2) Registra el motivo de la justificación. \newline
  3) El sistema verifica el contador de justificaciones del alumno (límite de
     3). \newline
  4) Crea la justificación, cambia el estado a \textit{Justificado} e
     incrementa el contador. \\
\addlinespace
Flujos alternativos &
  \textbf{3a. Límite alcanzado:} el alumno ya tiene 3 justificaciones
  $\rightarrow$ el sistema solicita confirmación del Director y genera una
  notificación interna dirigida al Director. \newline
  \textbf{3b. Confirmación del Director:} si el Director confirma, el ciclo de
  justificaciones se reinicia y se registra la nueva justificación. \\
\addlinespace
Postcondición &
  La asistencia queda en estado \textit{Justificado}, el contador del alumno se
  actualiza y, de corresponder, el Director recibe una notificación. \\
\bottomrule
\end{tabularx}
\end{table}

% ── CU04 ──────────────────────────────────────────────────────────
\subsection{CU04 -- Generar reporte de asistencia}

\begin{table}[H]
\centering
\begin{tabularx}{\linewidth}{C X}
\toprule
\textbf{Campo} & \textbf{Descripción} \\
\midrule
Actor & Director. \\
\addlinespace
Precondición & El usuario autenticado tiene rol Director y está activo. \\
\addlinespace
Flujo principal &
  1) El Director selecciona los filtros del reporte: fecha, sección y estado.
     \newline
  2) El sistema muestra una previsualización de las asistencias que coinciden.
     \newline
  3) El Director elige el formato de descarga: Excel (CSV) o PDF. \newline
  4) El sistema genera y entrega el archivo. \\
\addlinespace
Flujos alternativos &
  \textbf{1a. Sin filtros de fecha:} el sistema asume la fecha actual por
  defecto. \\
\addlinespace
Postcondición &
  El Director obtiene el archivo de reporte. El PDF incluye, además del estado
  de cada alumno, el nombre del apoderado y su teléfono de contacto. \\
\bottomrule
\end{tabularx}
\end{table}

% ── CU05 ──────────────────────────────────────────────────────────
\subsection{CU05 -- Crear un nuevo colegio (panel del dueño)}

\begin{table}[H]
\centering
\begin{tabularx}{\linewidth}{C X}
\toprule
\textbf{Campo} & \textbf{Descripción} \\
\midrule
Actor & Dueño de la plataforma (superusuario). \\
\addlinespace
Precondición &
  El dueño está autenticado en el panel de la plataforma; el subdominio
  elegido no está en uso. \\
\addlinespace
Flujo principal &
  1) El dueño completa el formulario de alta: nombre del colegio, subdominio,
     correo de contacto y contraseña del Director inicial. \newline
  2) El sistema crea un nuevo \textit{schema} aislado en PostgreSQL para el
     colegio y ejecuta sus migraciones. \newline
  3) Registra el dominio del colegio (subdominio + sufijo del dominio). \newline
  4) Crea, dentro del schema del colegio, el usuario \textit{Director} inicial.
     \newline
  5) Devuelve la URL de acceso del colegio. \\
\addlinespace
Flujos alternativos &
  \textbf{1a. Datos inválidos o subdominio duplicado:} el sistema rechaza el
  alta y muestra el error de validación. \\
\addlinespace
Postcondición &
  El colegio queda operativo con su subdominio activo bajo
  \texttt{yachayqr.com} y un Director listo para iniciar sesión. \\
\bottomrule
\end{tabularx}
\end{table}

% ── CU06 ──────────────────────────────────────────────────────────
\subsection{CU06 -- Gestionar alumnos y apoderados}

\begin{table}[H]
\centering
\begin{tabularx}{\linewidth}{C X}
\toprule
\textbf{Campo} & \textbf{Descripción} \\
\midrule
Actor & Director / Auxiliar (el Psicólogo solo consulta). \\
\addlinespace
Precondición & El usuario está autenticado, activo y con permiso de gestión. \\
\addlinespace
Flujo principal &
  1) El actor crea o edita un alumno: datos personales, DNI, foto, grado y
     sección. \newline
  2) Asigna al alumno un apoderado con su teléfono de WhatsApp y parentesco.
     \newline
  3) El sistema establece el código de barras del alumno igual a su DNI.
     \newline
  4) El actor puede generar el carnet del alumno en PDF (individual o de toda
     la sección). \\
\addlinespace
Flujos alternativos &
  \textbf{1a. DNI duplicado:} el sistema rechaza el registro por unicidad del
  DNI. \newline
  \textbf{4a. Sección sin alumnos activos:} al generar carnets por sección, si
  no hay alumnos el sistema informa el caso. \\
\addlinespace
Postcondición &
  El alumno queda registrado y asociado a un apoderado; su carnet con código de
  barras y QR está disponible para impresión. \\
\bottomrule
\end{tabularx}
\end{table}

% ════════════════════════════════════════════════════════════════
\section{Diagrama de Casos de Uso}
% ════════════════════════════════════════════════════════════════

La Figura~\ref{fig:cu} presenta el diagrama de casos de uso del sistema. Los
actores se ubican a los lados y los casos de uso, dentro del marco del sistema
\textbf{YachayQR}.

\begin{figure}[H]
\centering
\begin{tikzpicture}[
    every node/.style={transform shape},
    uc/.style={draw=navy, thick, ellipse, fill=gold!25,
               minimum width=3.4cm, minimum height=0.95cm,
               align=center, font=\small\color{navy}},
    assoc/.style={navy, thick}
  ]

  % ── Marco del sistema ──
  \draw[navy, very thick, rounded corners=8pt] (-2.9,-4.8) rectangle (2.9,5.3);
  \node[font=\large\bfseries, text=navy] at (0,4.85) {YachayQR};

  % ── Casos de uso (óvalos) ──
  \node[uc] (uc1) at (0,3.9)  {Registrar\\asistencia};
  \node[uc] (uc2) at (0,2.35) {Cerrar sesión /\\notificar WhatsApp};
  \node[uc] (uc3) at (0,0.8)  {Justificar\\inasistencia};
  \node[uc] (uc4) at (0,-0.75){Generar\\reporte};
  \node[uc] (uc6) at (0,-2.3) {Gestionar\\alumnos};
  \node[uc] (uc5) at (0,-3.85){Crear\\colegio};

  % ── Actores izquierda ──
  \stickman{-5.7}{4.0}{Escáner}
  \stickman{-5.7}{1.6}{Auxiliar}
  \stickman{-5.7}{-1.8}{Director}

  % ── Actores derecha ──
  \stickman{5.7}{0.8}{Psicólogo}
  \stickman{5.7}{-3.85}{Dueño\\Plataforma}

  % ── Asociaciones actores izquierda ──
  \draw[assoc] (-5.4,4.0)  -- (uc1.west);
  \draw[assoc] (-5.4,1.6)  -- (uc1.west);
  \draw[assoc] (-5.4,1.6)  -- (uc2.west);
  \draw[assoc] (-5.4,1.6)  -- (uc6.west);
  \draw[assoc] (-5.4,-1.8) -- (uc2.west);
  \draw[assoc] (-5.4,-1.8) -- (uc3.west);
  \draw[assoc] (-5.4,-1.8) -- (uc4.west);
  \draw[assoc] (-5.4,-1.8) -- (uc6.west);

  % ── Asociaciones actores derecha ──
  \draw[assoc] (5.4,0.8)   -- (uc3.east);
  \draw[assoc] (5.4,-3.85) -- (uc5.east);

\end{tikzpicture}
\caption{Diagrama de casos de uso del sistema YachayQR.}
\label{fig:cu}
\end{figure}

% ════════════════════════════════════════════════════════════════
\section{Stack Tecnológico}
% ════════════════════════════════════════════════════════════════

La Tabla~\ref{tab:stack} resume las tecnologías empleadas en cada capa del
sistema.

\begin{table}[H]
\centering
\begin{tabularx}{\linewidth}{p{3.2cm} X p{3.4cm}}
\toprule
\textbf{Capa} & \textbf{Tecnología} & \textbf{Versión} \\
\midrule
Backend &
  Django, Django REST Framework, django-tenants, Celery, Redis &
  Django 5.2 \\
\addlinespace
Frontend &
  React, Vite, React Router &
  React 19 / Router v7 \\
\addlinespace
Base de datos &
  PostgreSQL (un \textit{schema} por colegio) &
  PostgreSQL 16 \\
\addlinespace
Infraestructura &
  DigitalOcean, Nginx, Gunicorn, Let's Encrypt (SSL) &
  Ubuntu Server \\
\addlinespace
Servicios externos &
  Cloudflare (DNS + Turnstile anti-bot), Meta WhatsApp Cloud API &
  -- \\
\bottomrule
\end{tabularx}
\caption{Stack tecnológico de YachayQR.}
\label{tab:stack}
\end{table}

\subsection{Servicios externos: Cloudflare y Meta}

\textbf{Cloudflare} es la plataforma encargada de la gestión del dominio y la
seguridad de acceso del sistema. Cumple dos funciones:

\begin{itemize}[leftmargin=*]
  \item \textbf{DNS con comodín (\textit{wildcard}).} El dominio
  \texttt{yachayqr.com} se administra en Cloudflare mediante un registro
  comodín \texttt{*.yachayqr.com} que apunta al servidor. Gracias a esto, cada
  vez que se da de alta un colegio, su subdominio (por ejemplo
  \texttt{iestacoasa.yachayqr.com}) queda accesible de inmediato, sin necesidad
  de configurar manualmente un nuevo registro DNS.
  \item \textbf{Turnstile (anti-bot).} Es la alternativa de Cloudflare al
  CAPTCHA tradicional: verifica que quien intenta iniciar sesión sea una
  persona y no un programa automatizado, sin pedir al usuario resolver
  acertijos. La validación se realiza en el backend, lo que protege el
  formulario de acceso contra ataques de fuerza bruta.
\end{itemize}

\textbf{Meta WhatsApp Cloud API} es la interfaz oficial de Meta (la empresa
detrás de WhatsApp) que permite enviar mensajes de WhatsApp de forma
programática desde un número de empresa verificado, sin depender de un teléfono
físico. YachayQR la utiliza para \textbf{notificar a los apoderados}: al cerrar
la jornada, una tarea en segundo plano (Celery) envía el aviso de inasistencia
o tardanza al número registrado de cada apoderado. Este envío es parte del
\textit{Plan Premium}; los colegios que no lo tienen activo operan con
normalidad, pero el sistema omite el envío de mensajes.

\end{document}
